\section{Main Results}

The project should be read as a sequence rather than as a single Phase 1D experiment: benchmark foundation, BudgetRAG context policies, adaptive allocation, RLAIF labels, reward/preference construction, offline selector evaluation, retriever diversity, and KV-cache motivation.

All results below are offline logged-candidate results. They are not human preference learning, not online RL, not DPO/PPO/GRPO, and not runtime KV-cache pruning. Context reward remains a non-default candidate.

\subsection{Key Findings}

\begin{center}
\small
\begin{tabular}{>{\raggedright\arraybackslash}p{4.0cm} >{\raggedright\arraybackslash}p{7.0cm} >{\raggedright\arraybackslash}p{3.1cm}}
\toprule
Finding & Evidence & Interpretation \\
\midrule
Full context-level RLAIF labels are available. & 192/192 MiMo context labels were merged; 177 were clean usable labels; sufficiency rate was 0.591. & Context supervision is no longer only a skeleton or subset. \\
Retrieved contexts contain noise. & Mean selected chunks 1.276 versus mean irrelevant chunks 3.708. & Answer-level reward can hide weak evidence contexts. \\
Context reward changes the optimization target. & Penalty 0.25 changed 140/192 reward rows; preferences increased from 722 to 952; mean changed delta was -0.212. & Context reward is useful for calibration, but still non-default. \\
Pairwise RLAIF mostly agrees with scalar reward. & MiMo-50 pairwise audit agreement was 0.854 over non-ambiguous decisions; mean confidence was 0.914. & Scalar reward is not arbitrary, but needs calibration for near-tie answers. \\
Secondary judge supports the insufficiency signal. & DeepSeek v4 Flash agreed with MiMo on 80/83 comparable high-risk sufficiency decisions, agreement 0.964. & MiMo is not uniquely harsh on this targeted subset. \\
Retriever-diversity coverage is fixed. & Retrieval-only run produced 2250 actions across BM25, graph-BM25, and hybrid-RRF; A1-medium generation produced 1500/1500 non-empty rows with 0 errors and 0 empty answers. & Logged action coverage now supports retrieval-context allocation experiments. \\
A1 answer-level quality is now measured. & A1 has 1500/1500 valid MiMo V2.5 answer labels, 1460 clean usable labels, 1460 scored rewards, and graph-BM25 has the best mean answer-level reward among the three retrievers. & This is an answer-level signal, not a final retriever ranking. \\
A1 context-level evidence is now measured on a stratified subset. & The 600-row A1 context pass has 548 clean usable labels, sufficiency 0.6985, context quality 0.6431, evidence support 0.6567, and 3.487 irrelevant chunks on average. & Context supervision is no longer pending, but it covers 600/1500 rows. \\
A1 answer and context signals diverge. & Graph-BM25 is strongest at answer level; hybrid-RRF has the highest context sufficiency/support in the 600-row subset, while graph-BM25 has lower token cost and higher KV savings. & Retriever ranking is metric-dependent and should remain query-conditioned. \\
A1 selector signal is positive but costly. & With context reward 0.25, \code{linear\_reward\_model} remains the strongest non-oracle selector: reward 0.669 $\pm$ 0.081, quality 0.967 $\pm$ 0.026, oracle gap 0.145 $\pm$ 0.082. & The learned selector still uses high-cost actions; context reward remains non-default. \\
\bottomrule
\end{tabular}
\end{center}

\subsection{Phase Timeline and Deliverables}

{\scriptsize
\begin{longtable}{>{\raggedright\arraybackslash}p{1.5cm} >{\raggedright\arraybackslash}p{2.7cm} >{\raggedright\arraybackslash}p{4.8cm} >{\raggedright\arraybackslash}p{4.7cm}}
\toprule
Phase & Goal & Implementation & Key output/result \\
\midrule
Phase 1A & Benchmark foundation & Retriever registry, benchmark CLI, retrieval metrics, and traceable records. & A reusable output schema for later BudgetRAG and RLAIF stages. \\
Phase 1B & Fixed BudgetRAG policies & Legacy, character budget, per-document budget, score-density, sentence-trim, and evidence-aware policies. & 10-query SciFact smoke showed context compression and analytical KV savings can be measured independently from answer quality. \\
Phase 1B.1 & Traceable matrix outputs & Matrix runner and summarizer with stable raw output directories. & Phase outputs became auditable and joinable for later RLAIF builders. \\
Phase 1C & Adaptive heuristic & Deterministic selector using score gap, entropy, confidence, and document-dominance features. & Initial adaptive path was stable but conservative. \\
Phase 1C.1 & Larger retrieval validation & 50-query SciFact BM25 retrieval-only validation. & Adaptive heuristic selected \code{evidence-aware} for 48 queries and \code{per-doc-budget} for 2, but still selected budget 4000 for all 50. \\
Phase 1C.2 & Calibrated adaptive profiles & Conservative, balanced, and aggressive profiles with normalized gap/entropy thresholds. & Balanced/aggressive profiles created smaller-context stress tests; aggressive budget-4000 kept 2080 chars on average versus conservative 4000. \\
Phase 1C.3 & Generation feedback & Multi-model generation, RAGAS join, MiMo answer labels, and long-context MiMo runs. & 192 joined RLAIF action rows and 192 MiMo answer labels; MiMo quality vs RAGAS answer relevancy correlation 0.277. \\
Phase 1D & RLAIF reward/preference and offline selector & Answer/context/pairwise labelers, reward/preference builder, query-level split, selector baselines, and audits. & Full context labels, reward ablations, pairwise audit, action coverage diagnostics, and six-seed selector sweeps. \\
Retriever-diversity extension & Multi-retriever logged actions & BM25, graph-BM25, and hybrid-RRF retrieval-only and generation matrices. & Retrieval-only 2250-row coverage run; A1-medium 1500-row generation gate with complete answer labels and a 600-row stratified context-label pass. \\
\bottomrule
\end{longtable}
}

\subsection{Full Context-Level RLAIF Labels}

The first full MiMo context-label pass merged all 192 Phase 1D selector-smoke action rows. Validation found no missing action ids, no unknown action ids, no duplicate action ids, no duplicate conflicts, and no invalid JSON labels.

\begin{center}
\small
\begin{tabular}{l r}
\toprule
Metric & Value \\
\midrule
Merged context labels & 192 \\
Clean usable labels & 177 \\
Ambiguous labels & 15 \\
Invalid JSON labels & 0 \\
Sufficient contexts & 110 \\
Insufficient contexts & 76 \\
Sufficiency rate & 0.591 \\
Mean selected chunks & 1.276 \\
Mean redundant chunks & 0.286 \\
Mean irrelevant chunks & 3.708 \\
Mean context quality & 0.602 \\
Mean evidence support & 0.556 \\
Mean minimality & 0.947 \\
\bottomrule
\end{tabular}
\end{center}

The important evidence is not only the sufficiency rate. It is the gap between selected and irrelevant chunks: the judge typically identifies a small evidence subset while marking several retrieved chunks as irrelevant. This supports the BudgetRAG/MemAlign direction: context selection should be supervised at the chunk level, not only judged after answer generation.

\subsection{Context Reward Ablation}

The context labels were merged into reward construction only through explicit, non-default candidates. Ambiguous context labels fall back to answer-level feedback. The answer-only reward remains the baseline.

\begin{center}
\scriptsize
\begin{tabular}{>{\raggedright\arraybackslash}p{2.7cm} r r r r >{\raggedright\arraybackslash}p{3.0cm}}
\toprule
Setting & Preferences & Changed rows & Negative deltas & Mean changed delta & Interpretation \\
\midrule
answer-only & 722 & -- & -- & -- & Baseline reward from answer labels/fallback feedback. \\
context penalty 0.25 & 952 & 140 & 108 & -0.212 & Conservative context candidate; useful for further audit. \\
context penalty 0.50 & 946 & 140 & 108 & -0.316 & Stronger penalty; lower absolute reward scale. \\
context penalty 1.00 & 944 & 140 & 108 & -0.525 & Too harsh for default use; diagnostic only. \\
\bottomrule
\end{tabular}
\end{center}

The reward landscape changes materially: penalty 0.25 increases preference count from 722 to 952 and changes 140/192 reward rows. Most changes are negative because insufficient or noisy contexts are penalized. This is expected; it means the context labels are exposing evidence failures that answer-level labels alone do not represent.

\subsection{Pairwise RLAIF Audit}

Direct pairwise RLAIF was used to audit whether scalar reward-derived preferences align with an independent pairwise judge. In the MiMo-50 audit, Action A was the scalar reward-chosen action and Action B was the rejected action.

\begin{center}
\small
\begin{tabular}{l r}
\toprule
Metric & Value \\
\midrule
Pair labels & 50 \\
Valid JSON & 50 \\
Ambiguous/error rows & 2 \\
A wins & 41 \\
B wins & 7 \\
Agreement over non-ambiguous decisions & 0.854 \\
Mean confidence & 0.914 \\
Small quality/support delta pairs & 38 \\
Cheaper wins when quality/support tied & 35 \\
Scalar-over-quality disagreements & 5 \\
\bottomrule
\end{tabular}
\end{center}

The disagreement pattern is narrow: all seven direct-judge disagreements in the MiMo-50 audit came from \code{query\_id=128}. In those cases the scalar reward favored slightly higher quality/support scores, while the pairwise judge considered both answers acceptable or both correct abstentions and then preferred the cheaper action. The calibration diagnostic therefore suggests a candidate tie band of quality 0.10 and support 0.20, but this remains an audit-derived candidate, not a default reward rule.

\subsection{DeepSeek Secondary-Judge Audit}

The retriever-diverse 300-row subset was used to select 100 high-risk context rows: MiMo-insufficient contexts, large negative context-reward deltas, high answer quality with low context support, and rows with many irrelevant chunks. DeepSeek v4 Flash was then used as a secondary context judge.

\begin{center}
\small
\begin{tabular}{l r}
\toprule
Metric & Value \\
\midrule
Targeted rows & 100 \\
DeepSeek valid JSON labels & 100 \\
DeepSeek ambiguous labels & 12 \\
DeepSeek invalid JSON/errors & 0 \\
Comparable MiMo-vs-DeepSeek sufficiency rows & 83 \\
Agreement & 80/83 = 0.964 \\
Consensus-insufficient rows & 76 \\
High-disagreement rows & 3 \\
MiMo-harsh rows & 1 \\
Evidence-support Pearson correlation & 0.905 \\
Context-quality Pearson correlation & 0.476 \\
\bottomrule
\end{tabular}
\end{center}

This audit does not turn AI feedback into ground truth. It does, however, reduce the concern that MiMo is uniquely harsh on high-risk rows. The stronger agreement on evidence support than on raw context quality also suggests that sufficiency/support labels are safer reward inputs than an uncalibrated context-quality scalar.

\subsection{Offline Selector Results}

The selector is evaluated only over logged candidate actions. The artifact keeps \code{runtime\_default\_replacement=false}; it does not replace the runtime \code{adaptive-heuristic} policy.

\begin{center}
\small
\begin{tabular}{>{\raggedright\arraybackslash}p{4.4cm} r r r r}
\toprule
Policy & Reward & Quality & Coverage & Oracle gap \\
\midrule
cheapest & 0.577 $\pm$ 0.159 & 0.770 $\pm$ 0.152 & 1.000 $\pm$ 0.000 & 0.094 \\
best average & 0.618 $\pm$ 0.147 & 0.794 $\pm$ 0.125 & 0.911 $\pm$ 0.050 & 0.080 \\
shrinkage smoothed best average & 0.602 $\pm$ 0.117 & 0.773 $\pm$ 0.097 & 1.000 $\pm$ 0.000 & 0.070 \\
smoothed linear selector & 0.595 $\pm$ 0.137 & 0.767 $\pm$ 0.122 & 1.000 $\pm$ 0.000 & 0.076 \\
oracle logged & 0.671 $\pm$ 0.134 & 0.834 $\pm$ 0.138 & 1.000 $\pm$ 0.000 & 0.000 \\
\bottomrule
\end{tabular}
\end{center}

\textbf{Takeaway:} \code{shrinkage\_smoothed\_best\_average} is the best full-coverage non-oracle baseline by reward and oracle gap, while \code{best\_average} still has higher reward/quality on the query groups it covers. This indicates that action representation and logged-action coverage remain stronger bottlenecks than adding a more complex RL algorithm. Learned linear selectors are functional and full-coverage, but they do not yet robustly beat simple smoothed baselines.

\subsection{Retriever-Diversity Extension}

The first Phase 1D logs were effectively BM25-only. The retriever-diversity extension fixes this logged-action coverage gap, then starts adding quality supervision.

{\tiny
\setlength{\tabcolsep}{3pt}
\begin{longtable}{@{}>{\raggedright\arraybackslash}p{2.7cm} >{\raggedleft\arraybackslash}p{1.5cm} r >{\raggedleft\arraybackslash}p{1.4cm} >{\raggedright\arraybackslash}p{5.5cm}@{}}
\toprule
Run & Actions & Retrievers & Scored rewards & Purpose / conclusion \\
\midrule
Retrieval-only limit50 & 2250 & 3 & 0 & Coverage run only: every sampled query has BM25, graph-BM25, and hybrid-RRF rows across 45 action configurations. \\
300-row generation subset & 300 & 3 & 186 & Diagnostic quality run: graph-BM25 had slightly better answer support/context quality; hybrid-RRF was noisier, but \code{MAX\_COMPLETION\_TOKENS=256} caused 77 empty answers. \\
A1-medium generation + answer labels & 1500 & 3 & 1460 & Generation gate plus answer-level RLAIF: 1500/1500 valid MiMo V2.5 answer labels, 1460 clean usable labels, 0 invalid JSON, and 17026 answer-only preference pairs. \\
A1 stratified context labels + reward & 600 labels / 1500 rewards & 3 & 1492 & Context-level RLAIF subset: 548 clean usable context labels, sufficiency 0.6985, and a non-default context reward candidate changing 464 reward rows. \\
\bottomrule
\end{longtable}
}

\begin{center}
\small
\begin{tabular}{l r}
\toprule
A1 answer-label validation metric & Value \\
\midrule
Action rows & 1500 \\
Valid MiMo V2.5 labels & 1500 \\
Clean usable labels & 1460 \\
Ambiguous labels & 40 \\
Invalid JSON labels & 0 \\
Missing / unknown / duplicate action ids & 0 \\
Scored answer-only rewards & 1460 \\
Answer-only preferences & 17026 \\
\bottomrule
\end{tabular}
\end{center}

\begin{center}
\small
\begin{tabular}{l r r r r r r}
\toprule
Retriever & Rows & Quality & Support & Unsupported risk & Reward & KV cost \\
\midrule
BM25 & 500 & 0.899 & 0.899 & 0.096 & 0.617 & 0.605 \\
graph-BM25 & 500 & 0.911 & 0.909 & 0.087 & 0.642 & 0.559 \\
hybrid-RRF & 500 & 0.862 & 0.862 & 0.141 & 0.546 & 0.673 \\
\bottomrule
\end{tabular}
\end{center}

\begin{center}
\small
\begin{tabular}{l r r r r r}
\toprule
Retriever & Clean rows & Suff. & Context quality & Evidence support & Token cost \\
\midrule
BM25 & 182 & 0.720 & 0.647 & 0.685 & 687.247 \\
graph-BM25 & 180 & 0.706 & 0.642 & 0.659 & 579.606 \\
hybrid-RRF & 186 & 0.780 & 0.682 & 0.730 & 689.935 \\
\bottomrule
\end{tabular}
\end{center}

{\tiny
\begin{longtable}{>{\raggedright\arraybackslash}p{2.8cm} r r r r r r}
\toprule
A1 answer-only selector & Cov. & Reward & Quality & Token & KV & Gap \\
\midrule
cheapest & 0.900 $\pm$ 0.058 & 0.270 $\pm$ 0.115 & 0.702 $\pm$ 0.056 & 0.249 $\pm$ 0.002 & 0.248 $\pm$ 0.000 & 0.546 $\pm$ 0.116 \\
best avg. & 1.000 $\pm$ 0.000 & 0.661 $\pm$ 0.111 & 0.937 $\pm$ 0.070 & 0.554 $\pm$ 0.135 & 0.758 $\pm$ 0.260 & 0.154 $\pm$ 0.111 \\
shrinkage avg. & 1.000 $\pm$ 0.000 & 0.659 $\pm$ 0.112 & 0.938 $\pm$ 0.072 & 0.577 $\pm$ 0.137 & 0.808 $\pm$ 0.265 & 0.156 $\pm$ 0.111 \\
linear model & 1.000 $\pm$ 0.000 & 0.734 $\pm$ 0.032 & 0.987 $\pm$ 0.018 & 0.656 $\pm$ 0.014 & 0.994 $\pm$ 0.001 & 0.081 $\pm$ 0.032 \\
smoothed linear & 0.967 $\pm$ 0.047 & 0.657 $\pm$ 0.126 & 0.931 $\pm$ 0.076 & 0.480 $\pm$ 0.121 & 0.659 $\pm$ 0.240 & 0.158 $\pm$ 0.125 \\
oracle & 1.000 $\pm$ 0.000 & 0.815 $\pm$ 0.001 & 1.000 $\pm$ 0.000 & 0.264 $\pm$ 0.008 & 0.249 $\pm$ 0.000 & 0.000 $\pm$ 0.000 \\
\bottomrule
\end{longtable}
}

\begin{center}
\small
\begin{tabular}{l r}
\toprule
A1 context-label validation metric & Value \\
\midrule
Context labels & 600 \\
Valid JSON labels & 598 \\
Clean usable labels & 548 \\
Ambiguous labels & 52 \\
Invalid JSON labels & 2 \\
Missing / unknown / duplicate action ids & 0 \\
Sufficiency rate & 0.6985 \\
Mean selected chunks & 1.2300 \\
Mean irrelevant chunks & 3.4867 \\
Mean context quality & 0.6431 \\
Mean evidence support & 0.6567 \\
\bottomrule
\end{tabular}
\end{center}

{\tiny
\begin{longtable}{>{\raggedright\arraybackslash}p{2.8cm} r r r r r r}
\toprule
A1 context-candidate selector & Cov. & Reward & Quality & Token & KV & Gap \\
\midrule
cheapest & 1.000 $\pm$ 0.000 & 0.358 $\pm$ 0.104 & 0.793 $\pm$ 0.041 & 0.249 $\pm$ 0.002 & 0.248 $\pm$ 0.000 & 0.456 $\pm$ 0.105 \\
best avg. & 0.983 $\pm$ 0.037 & 0.613 $\pm$ 0.127 & 0.906 $\pm$ 0.061 & 0.321 $\pm$ 0.052 & 0.298 $\pm$ 0.082 & 0.201 $\pm$ 0.126 \\
shrinkage avg. & 0.983 $\pm$ 0.037 & 0.555 $\pm$ 0.096 & 0.885 $\pm$ 0.045 & 0.320 $\pm$ 0.053 & 0.310 $\pm$ 0.079 & 0.259 $\pm$ 0.096 \\
linear model & 1.000 $\pm$ 0.000 & 0.669 $\pm$ 0.081 & 0.967 $\pm$ 0.026 & 0.623 $\pm$ 0.017 & 0.982 $\pm$ 0.028 & 0.145 $\pm$ 0.082 \\
smoothed linear & 0.983 $\pm$ 0.037 & 0.619 $\pm$ 0.133 & 0.925 $\pm$ 0.048 & 0.379 $\pm$ 0.083 & 0.473 $\pm$ 0.189 & 0.195 $\pm$ 0.131 \\
oracle & 1.000 $\pm$ 0.000 & 0.814 $\pm$ 0.003 & 1.000 $\pm$ 0.000 & 0.266 $\pm$ 0.014 & 0.273 $\pm$ 0.035 & 0.000 $\pm$ 0.000 \\
\bottomrule
\end{longtable}
}

\begin{center}
\small
\begin{tabular}{l r r r}
\toprule
Context-candidate selector & BM25 & graph-BM25 & hybrid-RRF \\
\midrule
best avg. & 42 & 17 & 1 \\
shrinkage avg. & 31 & 27 & 2 \\
linear model & 8 & 51 & 1 \\
smoothed linear & 38 & 19 & 3 \\
oracle & 31 & 20 & 9 \\
\bottomrule
\end{tabular}
\end{center}

\textbf{Takeaway:} the retriever-diversity extension has fixed the logged-action coverage problem and now has answer-level and stratified context-level quality labels. Graph-BM25 is the strongest mean answer-level retriever in A1. The 600-row context labels do not simply confirm that ranking: hybrid-RRF has the best context sufficiency/support, while graph-BM25 has lower token cost and more KV savings. With context reward 0.25, \code{linear\_reward\_model} is still the strongest non-oracle selector by reward and oracle gap, but it favors graph-BM25 and high KV-cost actions. The remaining caveat is calibration and coverage:

\begin{quote}
Retriever diversity coverage, answer-level A1 evaluation, and a 600-row stratified context-label pass are complete. A final BM25 vs graph-BM25 vs hybrid-RRF ranking remains provisional because context labels cover 600/1500 actions and answer-level/context-level signals disagree.
\end{quote}
